\documentclass[12pt]{article}
\usepackage[utf8]{inputenc}
\usepackage[T1]{fontenc}
\usepackage{hyperref}
\title{Siberian Tatar Language Resources for K-12 Curriculum Development}
\author{Compiled by Claude Sonnet 5.5 (claude-sonnet-5-5) for Andr\'as Kornai}
\date{}
\begin{document}
\maketitle

\section{Introduction}
Siberian Tatar (ISO 639-3: sty; Glottocode: sibe1250) is a Turkic language of the Russian Federation. This report rests on a single web search and is deliberately short. Many results concerned standard Tatar (Tatarstan) rather than Siberian Tatar, and those are not listed.

\section{Primary Educational Resources}
No Siberian Tatar primers or textbooks turned up in the search. A source quoting the linguist Maxim Sagidullin says that textbooks still need to be developed and that a grammar is the minimum next step. That source also reports that 53 schools in the Tyumen region have an ethno-cultural component with about 2 hours a week of Tatar language and literature. It is unclear whether this is Siberian Tatar or standard Tatar, and the article is undated.
\begin{itemize}
\item Report on Tatar-language teaching in the Tyumen region and the lack of textbooks: \url{https://tatar-congress.org/en/?p=9195}
\end{itemize}

\section{Linguistic and Reference Materials}
\begin{itemize}
\item Article on J.~Giganov's 1804 Russian--Tatar dictionary (Tobolsk), described as an early dialectological record of Siberian Tatar vocabulary: \url{https://journals.tsu.ru/vestnik/en/&amp;journal_page=archive&amp;id=868&amp;article_id=49459}
\item Euralex 2016 paper, a lexicography-related paper surfaced by the search (its relevance to Siberian Tatar was not confirmed): \url{https://euralex.org/wp-content/themes/euralex/proceedings/Euralex%202016/euralex_2016_052_p475.pdf}
\end{itemize}
Other works named in search summaries, without a URL for the work itself: Tumasheva, \emph{Dictionary of the dialects of Siberian Tatars} (Kazan University Press, 1992), and a recent Russian--Siberian Tatar dictionary whose title and publisher could not be identified.

\section{Religious and Community Resources}
No religious or community-organisation resources specific to Siberian Tatar were found.

\section{Digital and Audio Resources}
No audio, video, corpus or software resources specific to Siberian Tatar were found.

\section{Curriculum Development Context}
Little verifiable material exists in the search results. A curriculum would have to start from dialect dictionaries and new grammar and textbook writing, and would need local expertise in the Tyumen region. Further targeted searching in Russian is recommended.
\end{document}